% =====================================================================
%  Anatomía de un LLM — Fase V — Módulo 21
%  Aprendizaje por refuerzo y preferencias
%  Compilar: latexmk -r ../latexmkrc M21-preferencias.tex
% =====================================================================
\documentclass[11pt,a4paper,twoside,openany]{book}
\usepackage{../estilo-apuntes}

\newcommand{\ApuntesCurso}{Anatomía de un LLM}
\newcommand{\ApuntesModulo}{Módulo 21}

% Notación propia del módulo
\newcommand{\Sset}{\mathcal{S}}
\newcommand{\Aset}{\mathcal{A}}
\newcommand{\KL}{D_{\mathrm{KL}}}
\newcommand{\piref}{\pi_{\mathrm{ref}}}
\newcommand{\piold}{\pi_{\mathrm{old}}}
\newcommand{\clip}{\operatorname{clip}}
\newcommand{\media}{\operatorname{media}}
\newcommand{\desv}{\operatorname{desv}}
\newcommand{\Var}{\operatorname{Var}}
\newcommand{\LRM}{\mathcal{L}_{\mathrm{RM}}}
\newcommand{\LDPO}{\mathcal{L}_{\mathrm{DPO}}}
\newcommand{\pref}{\succ}
\makeatletter\g@addto@macro\UrlBreaks{\do\-\do\/}\makeatother

\hypersetup{pdftitle={Módulo 21. Aprendizaje por refuerzo y preferencias (resumen)},pdfauthor={Anatomía de un LLM}}

\makeatletter
\newlabel{fig:bucle}{{1.1}{4}{El bucle agente-entorno en la generación de texto. La transición es determinista; la única señal externa es la recompensa final}{figure.1.1}{}}
\newlabel{fig:bucle@cref}{{[figure][1][1]1.1}{[1][3][]4}}
\newlabel{ej:arbol}{{1.11}{4}{Valores y ventajas en un árbol de respuestas}{ejemplo.11}{}}
\newlabel{ej:arbol@cref}{{[ejemplo][11][]1.11}{[1][4][]4}}
\newlabel{eje:1-chatbot}{{1.1}{4}{}{ejercicio.1.1}{}}
\newlabel{eje:1-chatbot@cref}{{[ejercicio][1][1]1.1}{[1][4][]4}}
\newlabel{eje:1-contar}{{1.2}{4}{}{ejercicio.1.2}{}}
\newlabel{eje:1-contar@cref}{{[ejercicio][2][1]1.2}{[1][4][]4}}
\newlabel{eje:1-retornos}{{1.3}{5}{}{ejercicio.1.3}{}}
\newlabel{eje:1-retornos@cref}{{[ejercicio][3][1]1.3}{[1][5][]5}}
\newlabel{eje:1-arbol}{{1.4}{5}{}{ejercicio.1.4}{}}
\newlabel{eje:1-arbol@cref}{{[ejercicio][4][1]1.4}{[1][5][]5}}
\newlabel{eje:1-cs285}{{1.5}{5}{}{ejercicio.1.5}{}}
\newlabel{eje:1-cs285@cref}{{[ejercicio][5][1]1.5}{[1][5][]5}}
\newlabel{eje:1-contraccion}{{1.6}{5}{}{ejercicio.1.6}{}}
\newlabel{eje:1-contraccion@cref}{{[ejercicio][6][1]1.6}{[1][5][]5}}
\newlabel{ej:reinforce}{{2.7}{7}{Un paso de REINFORCE en un bandido}{ejemplo.7}{}}
\newlabel{ej:reinforce@cref}{{[ejemplo][7][]2.7}{[1][7][]7}}
\newlabel{ej:baseline-optima}{{2.12}{9}{}{ejemplo.12}{}}
\newlabel{ej:baseline-optima@cref}{{[ejemplo][12][]2.12}{[1][9][]9}}
\newlabel{ej:gae}{{2.15}{9}{Cálculo de GAE}{ejemplo.15}{}}
\newlabel{ej:gae@cref}{{[ejemplo][15][]2.15}{[1][9][]9}}
\newlabel{eje:2-baseline}{{2.1}{10}{}{ejercicio.2.1}{}}
\newlabel{eje:2-baseline@cref}{{[ejercicio][1][2]2.1}{[1][10][]10}}
\newlabel{eje:2-bernoulli}{{2.2}{10}{}{ejercicio.2.2}{}}
\newlabel{eje:2-bernoulli@cref}{{[ejercicio][2][2]2.2}{[1][10][]10}}
\newlabel{eje:2-paso}{{2.3}{10}{}{ejercicio.2.3}{}}
\newlabel{eje:2-paso@cref}{{[ejercicio][3][2]2.3}{[1][10][]10}}
\newlabel{eje:2-cs285}{{2.4}{10}{}{ejercicio.2.4}{}}
\newlabel{eje:2-cs285@cref}{{[ejercicio][4][2]2.4}{[1][10][]10}}
\newlabel{eje:2-gae}{{2.5}{10}{}{ejercicio.2.5}{}}
\newlabel{eje:2-gae@cref}{{[ejercicio][5][2]2.5}{[1][10][]10}}
\newlabel{eje:2-optima}{{2.6}{10}{}{ejercicio.2.6}{}}
\newlabel{eje:2-optima@cref}{{[ejercicio][6][2]2.6}{[1][10][]10}}
\newlabel{eje:2-varianza}{{2.7}{10}{}{ejercicio.2.7}{}}
\newlabel{eje:2-varianza@cref}{{[ejercicio][7][2]2.7}{[1][10][]10}}
\newlabel{ej:kl}{{3.6}{13}{}{ejemplo.6}{}}
\newlabel{ej:kl@cref}{{[ejemplo][6][]3.6}{[1][13][]13}}
\newlabel{ej:beta}{{3.8}{13}{El papel de $\beta $}{ejemplo.8}{}}
\newlabel{ej:beta@cref}{{[ejemplo][8][]3.8}{[1][13][]13}}
\newlabel{eje:3-pesimista}{{3.1}{14}{}{ejercicio.3.1}{}}
\newlabel{eje:3-pesimista@cref}{{[ejercicio][1][3]3.1}{[1][14][]14}}
\newlabel{eje:3-ppo}{{3.2}{14}{}{ejercicio.3.2}{}}
\newlabel{eje:3-ppo@cref}{{[ejercicio][2][3]3.2}{[1][14][]14}}
\newlabel{eje:3-grpoclip}{{3.3}{14}{}{ejercicio.3.3}{}}
\newlabel{eje:3-grpoclip@cref}{{[ejercicio][3][3]3.3}{[1][14][]14}}
\newlabel{eje:3-beta}{{3.4}{14}{}{ejercicio.3.4}{}}
\newlabel{eje:3-beta@cref}{{[ejercicio][4][3]3.4}{[1][14][]14}}
\newlabel{eje:3-estimadores}{{3.5}{14}{}{ejercicio.3.5}{}}
\newlabel{eje:3-estimadores@cref}{{[ejercicio][5][3]3.5}{[1][14][]14}}
\newlabel{eje:3-limites}{{3.6}{15}{}{ejercicio.3.6}{}}
\newlabel{eje:3-limites@cref}{{[ejercicio][6][3]3.6}{[1][15][]15}}
\newlabel{ej:bt}{{4.5}{17}{}{ejemplo.5}{}}
\newlabel{ej:bt@cref}{{[ejemplo][5][]4.5}{[1][17][]17}}
\newlabel{fig:rlhf}{{4.1}{18}{Las tres etapas de RLHF según InstructGPT~\cite {instructgpt}. El modelo SFT inicializa la política, sirve de referencia congelada para la penalización KL y, en general, también inicializa el modelo de recompensa}{figure.4.1}{}}
\newlabel{fig:rlhf@cref}{{[figure][1][4]4.1}{[1][18][]18}}
\newlabel{eje:4-conceptos}{{4.1}{19}{}{ejercicio.4.1}{}}
\newlabel{eje:4-conceptos@cref}{{[ejercicio][1][4]4.1}{[1][19][]19}}
\newlabel{eje:4-bt}{{4.2}{19}{}{ejercicio.4.2}{}}
\newlabel{eje:4-bt@cref}{{[ejercicio][2][4]4.2}{[1][19][]19}}
\newlabel{eje:4-cs234}{{4.3}{19}{}{ejercicio.4.3}{}}
\newlabel{eje:4-cs234@cref}{{[ejercicio][3][4]4.3}{[1][19][]19}}
\newlabel{eje:4-ranking}{{4.4}{19}{}{ejercicio.4.4}{}}
\newlabel{eje:4-ranking@cref}{{[ejercicio][4][4]4.4}{[1][19][]19}}
\newlabel{eje:4-mv}{{4.5}{19}{}{ejercicio.4.5}{}}
\newlabel{eje:4-mv@cref}{{[ejercicio][5][4]4.5}{[1][19][]19}}
\newlabel{eje:4-convexidad}{{4.6}{19}{}{ejercicio.4.6}{}}
\newlabel{eje:4-convexidad@cref}{{[ejercicio][6][4]4.6}{[1][19][]19}}
\newlabel{ej:dpo}{{5.9}{22}{Un paso de DPO}{ejemplo.9}{}}
\newlabel{ej:dpo@cref}{{[ejemplo][9][]5.9}{[1][22][]22}}
\newlabel{eje:5-vf}{{5.1}{22}{}{ejercicio.5.1}{}}
\newlabel{eje:5-vf@cref}{{[ejercicio][1][5]5.1}{[1][22][]22}}
\newlabel{eje:5-paso}{{5.2}{23}{}{ejercicio.5.2}{}}
\newlabel{eje:5-paso@cref}{{[ejercicio][2][5]5.2}{[1][22][]23}}
\newlabel{eje:5-cs336}{{5.3}{23}{}{ejercicio.5.3}{}}
\newlabel{eje:5-cs336@cref}{{[ejercicio][3][5]5.3}{[1][23][]23}}
\newlabel{eje:5-tabular}{{5.4}{23}{}{ejercicio.5.4}{}}
\newlabel{eje:5-tabular@cref}{{[ejercicio][4][5]5.4}{[1][23][]23}}
\newlabel{eje:5-gradiente}{{5.5}{23}{}{ejercicio.5.5}{}}
\newlabel{eje:5-gradiente@cref}{{[ejercicio][5][5]5.5}{[1][23][]23}}
\newlabel{eje:5-invariancia}{{5.6}{23}{}{ejercicio.5.6}{}}
\newlabel{eje:5-invariancia@cref}{{[ejercicio][6][5]5.6}{[1][23][]23}}
\newlabel{fig:cai}{{6.1}{25}{Las dos fases de Constitutional AI~\cite {cai}. El modelo ajustado en la fase~1 genera los pares de la fase~2}{figure.6.1}{}}
\newlabel{fig:cai@cref}{{[figure][1][6]6.1}{[1][24][]25}}
\newlabel{ej:grpo}{{6.6}{26}{}{ejemplo.6}{}}
\newlabel{ej:grpo@cref}{{[ejemplo][6][]6.6}{[1][26][]26}}
\newlabel{eje:6-cai}{{6.1}{28}{}{ejercicio.6.1}{}}
\newlabel{eje:6-cai@cref}{{[ejercicio][1][6]6.1}{[1][27][]28}}
\newlabel{eje:6-grpo}{{6.2}{28}{}{ejercicio.6.2}{}}
\newlabel{eje:6-grpo@cref}{{[ejercicio][2][6]6.2}{[1][28][]28}}
\newlabel{eje:6-drgrpo}{{6.3}{28}{}{ejercicio.6.3}{}}
\newlabel{eje:6-drgrpo@cref}{{[ejercicio][3][6]6.3}{[1][28][]28}}
\newlabel{eje:6-longitud}{{6.4}{28}{}{ejercicio.6.4}{}}
\newlabel{eje:6-longitud@cref}{{[ejercicio][4][6]6.4}{[1][28][]28}}
\newlabel{eje:6-cero}{{6.5}{28}{}{ejercicio.6.5}{}}
\newlabel{eje:6-cero@cref}{{[ejercicio][5][6]6.5}{[1][28][]28}}
\newlabel{eje:6-gao}{{6.6}{28}{}{ejercicio.6.6}{}}
\newlabel{eje:6-gao@cref}{{[ejercicio][6][6]6.6}{[1][28][]28}}
\newlabel{eje:6-bon}{{6.7}{28}{}{ejercicio.6.7}{}}
\newlabel{eje:6-bon@cref}{{[ejercicio][7][6]6.7}{[1][28][]28}}
\newlabel{ap:soluciones}{{A}{29}{Soluciones de los ejercicios}{appendix.A}{}}
\newlabel{ap:soluciones@cref}{{[appendix][1][2147483647]A}{[1][29][]29}}
\newlabel{sol:eje:1-chatbot}{{A.1}{29}{Tema 1}{section*.3}{}}
\newlabel{sol:eje:1-chatbot@cref}{{[subappendix][1][2147483647,1]A.1}{[1][29][]29}}
\newlabel{sol:eje:1-contar}{{A.1}{29}{Tema 1}{section*.4}{}}
\newlabel{sol:eje:1-contar@cref}{{[subappendix][1][2147483647,1]A.1}{[1][29][]29}}
\newlabel{sol:eje:1-retornos}{{A.1}{30}{Tema 1}{section*.5}{}}
\newlabel{sol:eje:1-retornos@cref}{{[subappendix][1][2147483647,1]A.1}{[1][30][]30}}
\newlabel{sol:eje:1-arbol}{{A.1}{30}{Tema 1}{section*.6}{}}
\newlabel{sol:eje:1-arbol@cref}{{[subappendix][1][2147483647,1]A.1}{[1][30][]30}}
\newlabel{sol:eje:1-cs285}{{A.1}{31}{Tema 1}{section*.7}{}}
\newlabel{sol:eje:1-cs285@cref}{{[subappendix][1][2147483647,1]A.1}{[1][30][]31}}
\newlabel{sol:eje:1-contraccion}{{A.1}{31}{Tema 1}{section*.8}{}}
\newlabel{sol:eje:1-contraccion@cref}{{[subappendix][1][2147483647,1]A.1}{[1][31][]31}}
\newlabel{sol:eje:2-baseline}{{A.2}{32}{Tema 2}{section*.9}{}}
\newlabel{sol:eje:2-baseline@cref}{{[subappendix][2][2147483647,1]A.2}{[1][31][]32}}
\newlabel{sol:eje:2-bernoulli}{{A.2}{32}{Tema 2}{section*.10}{}}
\newlabel{sol:eje:2-bernoulli@cref}{{[subappendix][2][2147483647,1]A.2}{[1][32][]32}}
\newlabel{sol:eje:2-paso}{{A.2}{32}{Tema 2}{section*.11}{}}
\newlabel{sol:eje:2-paso@cref}{{[subappendix][2][2147483647,1]A.2}{[1][32][]32}}
\newlabel{sol:eje:2-cs285}{{A.2}{33}{Tema 2}{section*.12}{}}
\newlabel{sol:eje:2-cs285@cref}{{[subappendix][2][2147483647,1]A.2}{[1][33][]33}}
\newlabel{sol:eje:2-gae}{{A.2}{33}{Tema 2}{section*.13}{}}
\newlabel{sol:eje:2-gae@cref}{{[subappendix][2][2147483647,1]A.2}{[1][33][]33}}
\newlabel{sol:eje:2-optima}{{A.2}{33}{Tema 2}{section*.14}{}}
\newlabel{sol:eje:2-optima@cref}{{[subappendix][2][2147483647,1]A.2}{[1][33][]33}}
\newlabel{sol:eje:2-varianza}{{A.2}{34}{Tema 2}{section*.15}{}}
\newlabel{sol:eje:2-varianza@cref}{{[subappendix][2][2147483647,1]A.2}{[1][34][]34}}
\newlabel{sol:eje:3-pesimista}{{A.3}{35}{Tema 3}{section*.16}{}}
\newlabel{sol:eje:3-pesimista@cref}{{[subappendix][3][2147483647,1]A.3}{[1][34][]35}}
\newlabel{sol:eje:3-ppo}{{A.3}{35}{Tema 3}{section*.17}{}}
\newlabel{sol:eje:3-ppo@cref}{{[subappendix][3][2147483647,1]A.3}{[1][35][]35}}
\newlabel{sol:eje:3-grpoclip}{{A.3}{35}{Tema 3}{section*.18}{}}
\newlabel{sol:eje:3-grpoclip@cref}{{[subappendix][3][2147483647,1]A.3}{[1][35][]35}}
\newlabel{sol:eje:3-beta}{{A.3}{36}{Tema 3}{section*.19}{}}
\newlabel{sol:eje:3-beta@cref}{{[subappendix][3][2147483647,1]A.3}{[1][36][]36}}
\newlabel{sol:eje:3-estimadores}{{A.3}{36}{Tema 3}{section*.20}{}}
\newlabel{sol:eje:3-estimadores@cref}{{[subappendix][3][2147483647,1]A.3}{[1][36][]36}}
\newlabel{sol:eje:3-limites}{{A.3}{37}{Tema 3}{section*.21}{}}
\newlabel{sol:eje:3-limites@cref}{{[subappendix][3][2147483647,1]A.3}{[1][36][]37}}
\newlabel{sol:eje:4-conceptos}{{A.4}{37}{Tema 4}{section*.22}{}}
\newlabel{sol:eje:4-conceptos@cref}{{[subappendix][4][2147483647,1]A.4}{[1][37][]37}}
\newlabel{sol:eje:4-bt}{{A.4}{37}{Tema 4}{section*.23}{}}
\newlabel{sol:eje:4-bt@cref}{{[subappendix][4][2147483647,1]A.4}{[1][37][]37}}
\newlabel{sol:eje:4-cs234}{{A.4}{38}{Tema 4}{section*.24}{}}
\newlabel{sol:eje:4-cs234@cref}{{[subappendix][4][2147483647,1]A.4}{[1][38][]38}}
\newlabel{sol:eje:4-ranking}{{A.4}{38}{Tema 4}{section*.25}{}}
\newlabel{sol:eje:4-ranking@cref}{{[subappendix][4][2147483647,1]A.4}{[1][38][]38}}
\newlabel{sol:eje:4-mv}{{A.4}{38}{Tema 4}{section*.26}{}}
\newlabel{sol:eje:4-mv@cref}{{[subappendix][4][2147483647,1]A.4}{[1][38][]38}}
\newlabel{sol:eje:4-convexidad}{{A.4}{39}{Tema 4}{section*.27}{}}
\newlabel{sol:eje:4-convexidad@cref}{{[subappendix][4][2147483647,1]A.4}{[1][39][]39}}
\newlabel{sol:eje:5-vf}{{A.5}{39}{Tema 5}{section*.28}{}}
\newlabel{sol:eje:5-vf@cref}{{[subappendix][5][2147483647,1]A.5}{[1][39][]39}}
\newlabel{sol:eje:5-paso}{{A.5}{40}{Tema 5}{section*.29}{}}
\newlabel{sol:eje:5-paso@cref}{{[subappendix][5][2147483647,1]A.5}{[1][40][]40}}
\newlabel{sol:eje:5-cs336}{{A.5}{40}{Tema 5}{section*.30}{}}
\newlabel{sol:eje:5-cs336@cref}{{[subappendix][5][2147483647,1]A.5}{[1][40][]40}}
\newlabel{sol:eje:5-tabular}{{A.5}{40}{Tema 5}{section*.31}{}}
\newlabel{sol:eje:5-tabular@cref}{{[subappendix][5][2147483647,1]A.5}{[1][40][]40}}
\newlabel{sol:eje:5-gradiente}{{A.5}{41}{Tema 5}{section*.32}{}}
\newlabel{sol:eje:5-gradiente@cref}{{[subappendix][5][2147483647,1]A.5}{[1][41][]41}}
\newlabel{sol:eje:5-invariancia}{{A.5}{41}{Tema 5}{section*.33}{}}
\newlabel{sol:eje:5-invariancia@cref}{{[subappendix][5][2147483647,1]A.5}{[1][41][]41}}
\newlabel{sol:eje:6-cai}{{A.6}{42}{Tema 6}{section*.34}{}}
\newlabel{sol:eje:6-cai@cref}{{[subappendix][6][2147483647,1]A.6}{[1][42][]42}}
\newlabel{sol:eje:6-grpo}{{A.6}{42}{Tema 6}{section*.35}{}}
\newlabel{sol:eje:6-grpo@cref}{{[subappendix][6][2147483647,1]A.6}{[1][42][]42}}
\newlabel{sol:eje:6-drgrpo}{{A.6}{43}{Tema 6}{section*.36}{}}
\newlabel{sol:eje:6-drgrpo@cref}{{[subappendix][6][2147483647,1]A.6}{[1][42][]43}}
\newlabel{sol:eje:6-longitud}{{A.6}{43}{Tema 6}{section*.37}{}}
\newlabel{sol:eje:6-longitud@cref}{{[subappendix][6][2147483647,1]A.6}{[1][43][]43}}
\newlabel{sol:eje:6-cero}{{A.6}{43}{Tema 6}{section*.38}{}}
\newlabel{sol:eje:6-cero@cref}{{[subappendix][6][2147483647,1]A.6}{[1][43][]43}}
\newlabel{sol:eje:6-gao}{{A.6}{44}{Tema 6}{section*.39}{}}
\newlabel{sol:eje:6-gao@cref}{{[subappendix][6][2147483647,1]A.6}{[1][44][]44}}
\newlabel{sol:eje:6-bon}{{A.6}{44}{Tema 6}{section*.40}{}}
\newlabel{sol:eje:6-bon@cref}{{[subappendix][6][2147483647,1]A.6}{[1][44][]44}}
\makeatother

\begin{document}
\frontmatter
\begin{titlepage}
\centering
\vspace*{2.5cm}
{\large\scshape Anatomía de un LLM}\par
\vspace{0.4cm}
{\normalsize Fundamentos matemáticos de los modelos de lenguaje}\par
\vspace{2.2cm}
\rule{\textwidth}{0.8pt}\par\vspace{0.5cm}
{\small\scshape Fase V \textperiodcentered\ Alineamiento, RL, agentes e interpretabilidad}\par\vspace{0.3cm}
{\Huge\bfseries Módulo 21\par}\vspace{0.3cm}
{\LARGE Aprendizaje por refuerzo y preferencias\par}
\vspace{0.2cm}
{\large Gradiente de política, PPO, modelos de recompensa, DPO y GRPO\par}
\vspace{0.5cm}\rule{\textwidth}{0.4pt}\par
\vspace{1.6cm}
{\large Resumen del módulo}\par
\vfill
{\small Edición de octubre de 2026}\par
\end{titlepage}

\chapter*{Sobre este resumen}
\addcontentsline{toc}{chapter}{Sobre este resumen}
\markboth{Sobre este resumen}{}

Este documento reúne lo indispensable de los apuntes del Módulo 21: el texto expositivo, todas las definiciones y todos los resultados (teoremas, proposiciones, lemas y corolarios), las tablas y las observaciones de contenido. Se han suprimido las demostraciones, los ejemplos, las lecturas, las figuras, los ejercicios y sus soluciones. La numeración de temas, secciones, enunciados y ecuaciones es la misma que en los apuntes, de modo que cualquier referencia puede buscarse allí; las que remiten a elementos suprimidos (un ejemplo, una figura, una ecuación de una demostración) conservan su número en los apuntes.

\section*{Notación}
\begin{center}\small
\begin{tabularx}{\textwidth}{@{}l>{\raggedright\arraybackslash}X@{}}
\toprule
Símbolo & Significado \\
\midrule
$\Sset$, $\Aset$, $p(s'\mid s,a)$ & estados, acciones y probabilidades de transición de un MDP \\
$\gamma\in[0,1]$, $G_t$ & factor de descuento y retorno desde el instante $t$ \\
$\tau=(s_0,a_0,r_1,s_1,\dots)$ & trayectoria o episodio \\
$\pi_\theta(a\mid s)$ & política con parámetros $\theta$ \\
$V^{\pi}$, $Q^{\pi}$, $A^{\pi}=Q^{\pi}-V^{\pi}$ & funciones de valor de estado, de acción y ventaja \\
$J(\theta)$ & rendimiento esperado de la política $\pi_\theta$ \\
$x$, $y=y_{1:T}$, $\V$ & instrucción (\emph{prompt}), respuesta de $T$ tokens y vocabulario \\
$r(x,y)$, $r_\phi$ & recompensa de una respuesta y modelo de recompensa con parámetros $\phi$ \\
$\piref$, $\piold$ & política de referencia (congelada) y política que generó las muestras \\
$\rho=\pi_\theta/\piold$, $\varepsilon$ & cociente de importancia y parámetro de recorte de PPO \\
$\beta>0$, $Z(x)$ & intensidad de la penalización KL y función de partición \\
$\KL(p\,\|\,q)$ & divergencia de Kullback--Leibler, en nats \\
$y_w\pref y_l$ & la respuesta $y_w$ se prefiere a $y_l$ \\
$\sigma(z)=1/(1+e^{-z})$ & función sigmoide \\
$e_a$ & vector canónico (\emph{one-hot}) de la acción $a$ \\
\bottomrule
\end{tabularx}
\end{center}


\tableofcontents

\mainmatter
% =====================================================================
\chapter{Procesos de decisión de Markov y el LLM como política}
\label{tema:1}
% =====================================================================

Este tema introduce el lenguaje del aprendizaje por refuerzo: agente, entorno, recompensa, retorno, política y funciones de valor. La presentación sigue el capítulo~3 de Sutton y Barto~\cite{sb} y la clase~2 del curso de David Silver~\cite{silver}; la última sección traslada el formalismo a la generación de texto, como en la clase~17 de CS336~\cite{cs336}.

\section{La interfaz agente-entorno}

En cada instante $t=0,1,2,\dots$ un \emph{agente} observa el estado $S_t$ del \emph{entorno}, elige una acción $A_t$ y recibe, en el instante siguiente, una recompensa numérica $R_{t+1}$ y un nuevo estado $S_{t+1}$. La sucesión resultante
\[
S_0,\,A_0,\,R_1,\,S_1,\,A_1,\,R_2,\,S_2,\,\dots
\]
se llama \emph{trayectoria}. Cuando la interacción termina en un estado terminal tras un número finito de pasos, la trayectoria se llama \emph{episodio}.

\begin{definicion}[Proceso de decisión de Markov finito]\label{def:mdp}
Un \emph{proceso de decisión de Markov} (MDP) finito es una tupla $(\Sset,\Aset,p,\rho_0,\gamma)$ formada por un conjunto finito de estados $\Sset$, un conjunto finito de acciones $\Aset$, una \emph{dinámica}
\[
p(s',r\mid s,a)=P\bigl(S_{t+1}=s',\,R_{t+1}=r\mid S_t=s,\,A_t=a\bigr),
\]
que para cada par $(s,a)$ es una distribución sobre $\Sset\times\R$ con soporte finito, una distribución inicial $\rho_0$ sobre $\Sset$ y un \emph{factor de descuento} $\gamma\in[0,1]$. La palabra «Markov» expresa que la distribución de $(S_{t+1},R_{t+1})$ depende de la historia solo a través de $(S_t,A_t)$.
\end{definicion}

\begin{definicion}[Política]\label{def:politica}
Una \emph{política} (estacionaria y estocástica) es una aplicación que asigna a cada estado $s$ una distribución $\pi(\cdot\mid s)$ sobre $\Aset$. Una política \emph{parametrizada} $\pi_\theta$ depende de un vector de parámetros $\theta\in\R^{d}$.
\end{definicion}

Fijados un MDP y una política, la probabilidad de un prefijo de trayectoria $\tau=(s_0,a_0,r_1,s_1,\dots,s_T)$ se obtiene por la regla de la cadena (módulo~1):
\begin{equation}\label{eq:ptau}
p_\pi(\tau)=\rho_0(s_0)\prod_{t=0}^{T-1}\pi(a_t\mid s_t)\,p(s_{t+1},r_{t+1}\mid s_t,a_t).
\end{equation}
Esta factorización separa lo que controla el agente (los factores $\pi$) de lo que impone el entorno (los factores $p$), y es la base del teorema del gradiente de política del \cref{tema:2}.

\section{Recompensa y retorno}

El objetivo del agente no es maximizar la recompensa inmediata sino una cantidad acumulada~\cite[cap.~3]{sb}.

\begin{definicion}[Retorno]\label{def:retorno}
El \emph{retorno} desde el instante $t$ es
\[
G_t=R_{t+1}+\gamma R_{t+2}+\gamma^{2}R_{t+3}+\dots=\sum_{k=0}^{\infty}\gamma^{k}R_{t+k+1}.
\]
En una tarea episódica de horizonte $T$ se toma $R_k=0$ para $k>T$, de modo que la suma es finita y se admite $\gamma=1$.
\end{definicion}

\begin{proposicion}\label{prop:retorno}
\begin{enumerate}
  \item Para todo $t$, $G_t=R_{t+1}+\gamma\,G_{t+1}$.
  \item Si $\abs{R_k}\le R_{\max}$ para todo $k$ y $\gamma<1$, entonces $\abs{G_t}\le R_{\max}/(1-\gamma)$.
\end{enumerate}
\end{proposicion}

La recursión (i) permite calcular todos los retornos de un episodio en una sola pasada hacia atrás, empezando por $G_T=0$ (\cref{eje:1-retornos}).

\section{Funciones de valor y ecuación de Bellman}

\begin{definicion}[Funciones de valor y ventaja]\label{def:valor}
Sea $\pi$ una política. La \emph{función de valor de estado} y la \emph{función de valor de acción} son
\[
V^{\pi}(s)=\E_{\pi}\bigl[G_t\mid S_t=s\bigr],\qquad Q^{\pi}(s,a)=\E_{\pi}\bigl[G_t\mid S_t=s,\,A_t=a\bigr],
\]
donde $\E_\pi$ indica que las acciones futuras se eligen con $\pi$. La \emph{función ventaja} es $A^{\pi}(s,a)=Q^{\pi}(s,a)-V^{\pi}(s)$.
\end{definicion}

La ventaja mide cuánto mejor (o peor) que la media de la propia política es elegir $a$ en $s$ y continuar después con $\pi$. Es la cantidad que pondera los gradientes de todos los métodos de los temas~2, 3 y~6.

\begin{teorema}[Ecuación de Bellman para $V^{\pi}$]\label{teo:bellman}
Para todo estado $s$,
\[
V^{\pi}(s)=\sum_{a}\pi(a\mid s)\,Q^{\pi}(s,a),\qquad Q^{\pi}(s,a)=\sum_{s',r}p(s',r\mid s,a)\bigl[r+\gamma V^{\pi}(s')\bigr].
\]
En particular, $V^{\pi}(s)=\sum_{a}\pi(a\mid s)\sum_{s',r}p(s',r\mid s,a)\bigl[r+\gamma V^{\pi}(s')\bigr]$.
\end{teorema}

\begin{corolario}\label{cor:ventaja-media}
Para todo estado $s$, $\displaystyle\sum_{a}\pi(a\mid s)\,A^{\pi}(s,a)=0$.
\end{corolario}

La ecuación de Bellman es un sistema lineal de $\abs{\Sset}$ ecuaciones con $\abs{\Sset}$ incógnitas. Para $\gamma<1$ tiene solución única (\cref{eje:1-contraccion}); el capítulo~4 de~\cite{sb} la resuelve de forma iterativa.

\section{El modelo de lenguaje como política}

La generación autorregresiva encaja en la \cref{def:mdp} con una dinámica especialmente sencilla~\cite{cs336}.

\begin{definicion}[MDP de tokens]\label{def:mdp-tokens}
Sean $\V$ un vocabulario con un símbolo de fin y $x$ una instrucción. El \emph{MDP de tokens} asociado a $x$ tiene como estados las secuencias $s_t=(x,y_1,\dots,y_t)$, como acciones los tokens $a_t=y_{t+1}\in\V$, una transición determinista $s_{t+1}=(s_t,a_t)$ (concatenación), estado inicial $s_0=x$ y estados terminales los que acaban en el símbolo de fin o alcanzan una longitud máxima. La recompensa es nula salvo en la transición final, en la que vale $r(x,y)$ para la respuesta completa $y=y_{1:T}$, y se toma $\gamma=1$.
\end{definicion}

\begin{proposicion}\label{prop:mdp-tokens}
En el MDP de tokens, con una política $\pi_\theta(y_{t+1}\mid x,y_{1:t})$:
\begin{enumerate}
  \item la probabilidad de una trayectoria coincide con la de la respuesta, $p_{\pi_\theta}(\tau)=\pi_\theta(y\mid x)=\prod_{t=1}^{T}\pi_\theta(y_t\mid x,y_{<t})$;
  \item el retorno desde cualquier instante es $G_t=r(x,y)$, y $V^{\pi_\theta}(x)=\E_{y\sim\pi_\theta(\cdot\mid x)}[r(x,y)]$;
  \item para un prefijo $s_t$ y un token $a$, $Q^{\pi_\theta}(s_t,a)=V^{\pi_\theta}(s_t a)$, donde $s_t a$ denota la concatenación.
\end{enumerate}
\end{proposicion}

\begin{table}[h]
\centering\small
\begin{tabular}{@{}ll@{}}
\toprule
Elemento del MDP & En un modelo de lenguaje \\
\midrule
Estado $s_t$ & instrucción $x$ más los tokens generados $y_{1:t}$ \\
Acción $a_t$ & siguiente token $y_{t+1}$; $\abs{\Aset}=\abs{\V}$, del orden de $10^{4}$--$10^{5}$ \\
Transición & determinista y conocida: concatenar el token \\
Política $\pi_\theta(a\mid s)$ & softmax de los logits del Transformer (módulo~14) \\
Recompensa & normalmente solo al final: $r(x,y)$ \\
Episodio & una respuesta completa, hasta el símbolo de fin \\
\bottomrule
\end{tabular}
\caption{Correspondencia entre los elementos de un MDP y la generación de texto.}
\label{tab:mdp-llm}
\end{table}

\addtocounter{figure}{1}%

\begin{observacion}[Bandido contextual]\label{obs:bandido}
Como la transición es conocida y la recompensa llega solo al final, la respuesta entera $y$ puede verse como una única acción tomada en el contexto $x$ con probabilidad $\pi_\theta(y\mid x)$. Este caso particular de MDP de un solo paso se llama \emph{bandido contextual}. Los métodos de preferencias de los temas~4 y~5 trabajan a este nivel; PPO (\cref{tema:3}) trabaja token a token con ventajas $A^{\pi}(s_t,a_t)$. El número de acciones del bandido es enorme: con $\abs{\V}=5\cdot 10^{4}$ y $T=100$ hay $\abs{\V}^{T}\approx 10^{470}$ respuestas, por lo que ninguna suma sobre $y$ puede calcularse explícitamente.
\end{observacion}

\addtocounter{definicion}{1}%


% =====================================================================
\chapter{Gradiente de política}
\label{tema:2}
% =====================================================================

Este tema deduce el gradiente del rendimiento esperado respecto a los parámetros de la política y lo convierte en un algoritmo estimable por muestreo, REINFORCE~\cite{williams}. Después estudia la varianza del estimador y las dos técnicas que la reducen: la línea base y la ventaja estimada por un crítico, incluida la estimación generalizada de la ventaja (GAE)~\cite{gae}. Se siguen las §13.1--13.5 de Sutton y Barto~\cite{sb} y las clases~5 y~6 de CS285~\cite{cs285}.

\section{Políticas parametrizadas y función de puntuación}

\begin{definicion}[Rendimiento y función de puntuación]\label{def:score}
Sea $\pi_\theta$ una política parametrizada en un MDP episódico. El \emph{rendimiento} es $J(\theta)=\E_{\tau\sim p_\theta}[R(\tau)]$, con $R(\tau)=\sum_{t=0}^{T-1}\gamma^{t}r_{t+1}$ y $p_\theta=p_{\pi_\theta}$ dada por \eqref{eq:ptau}. Si $\pi_\theta(a\mid s)>0$, el vector $\nabla_\theta\log\pi_\theta(a\mid s)$ se llama \emph{función de puntuación} (\emph{score}) o vector de elegibilidad.
\end{definicion}

La parametrización más común en espacios de acciones discretos es la softmax de unas preferencias o logits~\cite[§13.1]{sb}. En un LLM los logits son la salida del Transformer y dependen de $\theta$ de forma no lineal; el caso tabular, en el que los propios logits son los parámetros, permite calcular todo a mano.

\begin{proposicion}[Puntuación de la softmax]\label{prop:score-softmax}
Sea $\pi_\theta=\softmax(\theta)$, con $\theta\in\R^{K}$ y $\pi_\theta(a)=e^{\theta_a}/\sum_b e^{\theta_b}$. Entonces
\[
\nabla_\theta\log\pi_\theta(a)=e_a-\pi_\theta .
\]
Si los logits son $z=f_\theta(s)$, por la regla de la cadena $\nabla_\theta\log\pi_\theta(a\mid s)=J_f(s)^{\top}(e_a-\pi_\theta(\cdot\mid s))$, con $J_f$ la matriz jacobiana de $f_\theta$.
\end{proposicion}

El vector $e_a-\pi_\theta$ es exactamente el gradiente, cambiado de signo, de la entropía cruzada con etiqueta $a$ (módulo~8). Esta coincidencia explica la relación entre gradiente de política y ajuste supervisado que se establece más abajo.

\section{El teorema del gradiente de política}

\begin{lema}[Truco de la derivada logarítmica]\label{lem:logder}
Si $p_\theta(z)>0$ es diferenciable en $\theta$, entonces $\nabla_\theta p_\theta(z)=p_\theta(z)\,\nabla_\theta\log p_\theta(z)$.
\end{lema}

\begin{teorema}[Gradiente de política, forma de trayectorias]\label{teo:pg}
En un MDP episódico finito con política diferenciable y estrictamente positiva,
\[
\nabla_\theta J(\theta)=\E_{\tau\sim p_\theta}\Bigl[R(\tau)\sum_{t=0}^{T-1}\nabla_\theta\log\pi_\theta(a_t\mid s_t)\Bigr].
\]
\end{teorema}

Lo notable del teorema es que la dinámica $p$ desaparece del gradiente: basta con poder \emph{muestrear} trayectorias y evaluar su recompensa, aunque ni el entorno ni la recompensa sean derivables. En la forma de Sutton y Barto~\cite[§13.2]{sb}, el teorema se enuncia con la distribución de estados de la política y $Q^{\pi}$; la versión anterior es la de la clase~5 de CS285~\cite{cs285}.

\begin{corolario}[Gradiente de política en un LLM]\label{cor:pg-llm}
En el MDP de tokens (\cref{def:mdp-tokens}),
\[
\nabla_\theta J(\theta)=\E_{x}\,\E_{y\sim\pi_\theta(\cdot\mid x)}\Bigl[r(x,y)\sum_{t=1}^{T}\nabla_\theta\log\pi_\theta(y_t\mid x,y_{<t})\Bigr].
\]
\end{corolario}

El sumatorio $\sum_t\nabla_\theta\log\pi_\theta(y_t\mid x,y_{<t})$ es el gradiente de la log-verosimilitud de $y$, el mismo que maximiza el SFT del módulo~20. El gradiente de política es, por tanto, un SFT sobre las propias muestras del modelo, ponderado por su recompensa~\cite{cs336}.

\section{REINFORCE}

\begin{definicion}[Estimador REINFORCE]\label{def:reinforce}
Dadas $N$ trayectorias independientes $\tau^{(1)},\dots,\tau^{(N)}\sim p_\theta$, el estimador REINFORCE del gradiente es
\[
\widehat{g}=\frac1N\sum_{i=1}^{N}R(\tau^{(i)})\sum_{t}\nabla_\theta\log\pi_\theta(a^{(i)}_t\mid s^{(i)}_t),
\]
y el algoritmo actualiza $\theta\leftarrow\theta+\eta\,\widehat{g}$ con una tasa de aprendizaje $\eta>0$.
\end{definicion}

Por el \cref{teo:pg}, $\E[\widehat g]=\nabla_\theta J$: el estimador es insesgado. Williams~\cite{williams} lo introdujo en 1992 como un caso de una familia de algoritmos que siguen el gradiente en esperanza.

\addtocounter{definicion}{1}%

\section{Causalidad y recompensa futura}

\begin{proposicion}[Las recompensas pasadas no aportan]\label{prop:causalidad}
Si $t'<t$, entonces $\E_{p_\theta}\bigl[r_{t'+1}\,\nabla_\theta\log\pi_\theta(a_t\mid s_t)\bigr]=0$. En consecuencia,
\[
\nabla_\theta J(\theta)=\E_{p_\theta}\Bigl[\sum_{t}\nabla_\theta\log\pi_\theta(a_t\mid s_t)\,\widehat{Q}_t\Bigr],\qquad \widehat Q_t=\sum_{t'\ge t}\gamma^{t'}r_{t'+1}.
\]
\end{proposicion}

La cantidad $\widehat Q_t$ (la «recompensa futura», \emph{reward-to-go}) es una estimación de una muestra de $Q^{\pi}(s_t,a_t)$, salvo el factor $\gamma^{t}$~\cite{cs285hw2}. En el MDP de tokens, con una sola recompensa final, todas las $\widehat Q_t$ valen $r(x,y)$ y la proposición no aporta nada; su utilidad aparece con recompensas intermedias, como la penalización KL por token del \cref{tema:3}.

\section{Línea base y varianza}

\begin{lema}[La puntuación tiene media nula]\label{lem:score-media}
Para todo estado $s$, $\E_{a\sim\pi_\theta(\cdot\mid s)}\bigl[\nabla_\theta\log\pi_\theta(a\mid s)\bigr]=0$.
\end{lema}

\begin{proposicion}[Línea base]\label{prop:baseline}
Sea $b(s)$ cualquier función del estado que no dependa de la acción. Entonces
\[
\nabla_\theta J(\theta)=\E_{p_\theta}\Bigl[\sum_t\nabla_\theta\log\pi_\theta(a_t\mid s_t)\bigl(\widehat Q_t-b(s_t)\bigr)\Bigr].
\]
\end{proposicion}

La línea base no cambia la esperanza, pero sí la varianza. En el \cref{ej:reinforce}, la varianza total (traza de la matriz de covarianzas) del estimador de una muestra $(r(a)-b)(e_a-\pi)$ vale $0{,}889$ con $b=0$ y $0{,}222$ con $b=J=1$. Si se suma~10 a todas las recompensas, con $b=0$ sube a $80{,}9$, mientras que con $b=J$ sigue siendo $0{,}222$: la línea base hace el método insensible al origen arbitrario de la escala de recompensas. La clase~5 de CS285 observa que la recompensa media no es la línea base óptima, aunque sí una muy buena~\cite{cs285}; la proposición siguiente da la óptima.

\begin{proposicion}[Línea base de varianza mínima]\label{prop:baseline-optima}
En un bandido con $g(a)=\nabla_\theta\log\pi_\theta(a)$, la varianza total del estimador $(r(a)-b)\,g(a)$ con $a\sim\pi_\theta$ es mínima para
\[
b^{*}=\frac{\E\bigl[r(a)\,\norm{g(a)}^{2}\bigr]}{\E\bigl[\norm{g(a)}^{2}\bigr]}.
\]
\end{proposicion}

\addtocounter{definicion}{1}%

\section{Ventaja, crítico y estimación generalizada de la ventaja}

La línea base natural en un MDP es $b(s_t)=V^{\pi}(s_t)$. Entonces $\widehat Q_t-V^{\pi}(s_t)$ es una estimación de la ventaja $A^{\pi}(s_t,a_t)$ de la \cref{def:valor}, y el gradiente queda $\E[\sum_t\nabla\log\pi_\theta(a_t\mid s_t)\,A^{\pi}(s_t,a_t)]$, salvo los factores de descuento. Como $V^{\pi}$ es desconocida, se aproxima con una red $V_\phi$, el \emph{crítico}, entrenada por regresión sobre los retornos observados. La política es el \emph{actor}: de ahí el nombre de métodos actor-crítico~\cite[§13.5]{sb}.

\begin{definicion}[Error de diferencia temporal y GAE]\label{def:gae}
Dado un crítico $V$, el \emph{error de diferencia temporal} es $\delta_t=r_{t+1}+\gamma V(s_{t+1})-V(s_t)$, con $V(s_T)=0$ en el estado terminal. Para $\lambda\in[0,1]$, la \emph{estimación generalizada de la ventaja} es~\cite{gae}
\[
\widehat A^{\mathrm{GAE}(\gamma,\lambda)}_t=\sum_{l=0}^{T-t-1}(\gamma\lambda)^{l}\,\delta_{t+l},
\]
que satisface la recursión $\widehat A_t=\delta_t+\gamma\lambda\,\widehat A_{t+1}$, con $\widehat A_T=0$.
\end{definicion}

\begin{proposicion}[Casos extremos de GAE]\label{prop:gae}
Con $\lambda=0$, $\widehat A_t=\delta_t$. Con $\lambda=1$, $\widehat A_t=\sum_{l\ge0}\gamma^{l}r_{t+l+1}-V(s_t)$, el retorno observado menos la línea base. Si $V=V^{\pi}$, entonces $\E[\delta_t\mid s_t,a_t]=A^{\pi}(s_t,a_t)$.
\end{proposicion}

El parámetro $\lambda$ regula el compromiso entre sesgo y varianza: con $\lambda=0$ se usa solo un paso de recompensa real y se confía en el crítico (poca varianza, sesgo si $V$ es inexacto); con $\lambda=1$ se usa el retorno completo (sin sesgo, mucha varianza)~\cite{cs285}. PPO calcula sus ventajas con GAE.

\addtocounter{definicion}{1}%


% =====================================================================
\chapter{Optimización proximal y regularización KL}
\label{tema:3}
% =====================================================================

REINFORCE usa cada lote de muestras para un único paso de gradiente. Este tema estudia cómo reutilizar las muestras sin alejarse demasiado de la política que las generó, mediante el muestreo de importancia y el recorte de PPO~\cite{ppo}, y cómo se impide que la política se aleje del modelo de partida mediante una penalización de Kullback--Leibler. La parte de PPO sigue la clase~9 de CS285~\cite{cs285}; la forma cerrada de la política óptima procede del apéndice~A.1 de Rafailov \emph{et al.}~\cite{dpo}, y el objetivo con KL, de la §3.5 de InstructGPT~\cite{instructgpt}.

\section{Muestreo de importancia y objetivo sustituto}

\begin{proposicion}[Muestreo de importancia]\label{prop:is}
Sean $p$ y $q$ distribuciones sobre un conjunto finito con $q(z)>0$ siempre que $p(z)>0$. Para toda función $f$,
\[
\E_{z\sim p}[f(z)]=\E_{z\sim q}\Bigl[\frac{p(z)}{q(z)}\,f(z)\Bigr].
\]
\end{proposicion}

Sea $\piold$ la política con la que se generaron las muestras. Para una pareja estado-acción se define el \emph{cociente de importancia}
\[
\rho_t(\theta)=\frac{\pi_\theta(a_t\mid s_t)}{\piold(a_t\mid s_t)}.
\]
El \emph{objetivo sustituto} $L(\theta)=\E_{\piold}[\rho_t(\theta)\,\widehat A_t]$ estima, con muestras de $\piold$, la mejora local de la política $\pi_\theta$; la aproximación consiste en ignorar que la distribución de estados también cambia con $\theta$, lo que es aceptable si $\pi_\theta$ y $\piold$ están próximas~\cite{cs285}.

\begin{proposicion}\label{prop:sustituto}
$\nabla_\theta\rho_t(\theta)=\rho_t(\theta)\,\nabla_\theta\log\pi_\theta(a_t\mid s_t)$. En particular, en $\theta=\theta_{\mathrm{old}}$, $\nabla_\theta L=\E_{\piold}[\nabla_\theta\log\pi_\theta(a_t\mid s_t)\,\widehat A_t]$, el gradiente de política con ventaja del \cref{tema:2}.
\end{proposicion}

Optimizar $L$ durante varios pasos con las mismas muestras es tentador, pero si $\rho_t$ se aleja mucho de~1 la estimación pierde validez y su varianza crece. TRPO resuelve el problema con una restricción explícita sobre la divergencia KL entre $\piold$ y $\pi_\theta$; PPO lo resuelve con un recorte, mucho más sencillo de implementar~\cite{cs285,ppo}.

\section{El objetivo recortado de PPO}

\begin{definicion}[Objetivo recortado]\label{def:ppo}
Sea $\varepsilon>0$ (típicamente $\varepsilon=0{,}2$) y sea $\clip(\rho,1-\varepsilon,1+\varepsilon)=\min\{\max\{\rho,1-\varepsilon\},1+\varepsilon\}$. El objetivo de PPO es~\cite[§3]{ppo}
\[
L^{\mathrm{CLIP}}(\theta)=\E_t\Bigl[\min\bigl(\rho_t(\theta)\widehat A_t,\ \clip(\rho_t(\theta),1-\varepsilon,1+\varepsilon)\,\widehat A_t\bigr)\Bigr].
\]
\end{definicion}

\begin{proposicion}[Análisis por casos]\label{prop:ppo-casos}
Para un término con ventaja $A$ y cociente $\rho$, sea $\ell(\rho)=\min(\rho A,\clip(\rho)A)$.
\begin{enumerate}
  \item $\ell(\rho)\le\rho A$ siempre: el recorte nunca aumenta el objetivo.
  \item Si $A>0$, $\ell(\rho)=\min(\rho,1+\varepsilon)\,A$: el término deja de crecer, y su derivada respecto a $\rho$ se anula, cuando $\rho>1+\varepsilon$.
  \item Si $A<0$, $\ell(\rho)=\max(\rho,1-\varepsilon)\,A$: la derivada se anula cuando $\rho<1-\varepsilon$.
\end{enumerate}
\end{proposicion}

\begin{table}[h]
\centering\small
\begin{tabular}{@{}rrrrrl@{}}
\toprule
$\rho$ & $A$ & $\rho A$ & $\clip(\rho)A$ & $\ell$ & Lectura \\
\midrule
$1{,}5$ & $+2$ & $3{,}0$ & $2{,}4$ & $2{,}4$ & acción buena que ya subió: sin incentivo para subir más \\
$0{,}5$ & $+2$ & $1{,}0$ & $1{,}6$ & $1{,}0$ & acción buena que bajó: el gradiente la corrige \\
$1{,}5$ & $-1$ & $-1{,}5$ & $-1{,}2$ & $-1{,}5$ & acción mala que subió: se penaliza entera \\
$0{,}5$ & $-1$ & $-0{,}5$ & $-0{,}8$ & $-0{,}8$ & acción mala que ya bajó: gradiente nulo \\
\bottomrule
\end{tabular}
\caption{Los cuatro casos del objetivo recortado con $\varepsilon=0{,}2$.}
\label{tab:ppo}
\end{table}

El mínimo convierte $L^{\mathrm{CLIP}}$ en una cota \emph{pesimista} del objetivo sustituto: el recorte elimina el incentivo para alejarse más de $\varepsilon$ en la dirección favorable, pero nunca impide corregir un movimiento desfavorable (\cref{tab:ppo}). En la práctica, PPO combina $L^{\mathrm{CLIP}}$ con una pérdida cuadrática para el crítico y una bonificación de entropía cuando ambos comparten parámetros, y el artículo describe también una variante con penalización KL adaptativa respecto a $\piold$~\cite[§4]{ppo}.

\section{Penalización KL respecto a un modelo de referencia}

En el ajuste por refuerzo de modelos de lenguaje no se maximiza solo la recompensa. Sea $\piref$ el modelo de partida congelado (el modelo SFT). El objetivo regularizado es~\cite{instructgpt}
\begin{equation}\label{eq:obj-kl}
\begin{aligned}&\max_{\pi}\ \E_{x}\bigl[J_\beta(\pi;x)\bigr],\\ &J_\beta(\pi;x)=\E_{y\sim\pi(\cdot\mid x)}[r(x,y)]-\beta\,\KL\bigl(\pi(\cdot\mid x)\,\|\,\piref(\cdot\mid x)\bigr),\end{aligned}
\end{equation}
con $\beta>0$.

La penalización cumple tres funciones: impide que la política explote errores del modelo de recompensa, que solo es fiable cerca de los datos con que se entrenó (\cref{tema:6}); conserva la fluidez y el conocimiento del preentrenamiento, y mantiene la diversidad, ya que sin ella el óptimo es una política determinista que solo emite la respuesta de máxima recompensa.

\begin{proposicion}[KL de secuencias como suma por tokens]\label{prop:kl-tokens}
Sean $\pi$ y $\piref$ modelos autorregresivos sobre respuestas. Entonces
\[
\begin{aligned}
\KL\bigl(\pi(\cdot\mid x)\,\|\,\piref(\cdot\mid x)\bigr)&=\E_{y\sim\pi}\Bigl[\sum_{t=1}^{T}\log\frac{\pi(y_t\mid x,y_{<t})}{\piref(y_t\mid x,y_{<t})}\Bigr]\\
&=\E_{y\sim\pi}\Bigl[\sum_{t=1}^{T}\KL\bigl(\pi(\cdot\mid s_{t-1})\,\|\,\piref(\cdot\mid s_{t-1})\bigr)\Bigr],
\end{aligned}
\]
con $s_{t-1}=(x,y_{<t})$.
\end{proposicion}

La proposición justifica la implementación habitual: a la recompensa de cada token muestreado se le resta $\beta\log\bigl(\pi_\theta(y_t\mid s_{t-1})/\piref(y_t\mid s_{t-1})\bigr)$, cuyo valor esperado acumulado es la KL. Cada término individual puede ser negativo; solo su esperanza es no negativa (\cref{eje:3-estimadores}).

\addtocounter{definicion}{1}%

\section{La política óptima en forma cerrada}

\begin{teorema}[Política óptima con penalización KL]\label{teo:optimo-kl}
Para cada $x$, el objetivo \eqref{eq:obj-kl} sobre las distribuciones $\pi(\cdot\mid x)$ alcanza su máximo, de forma única, en
\[
\pi^{*}(y\mid x)=\frac{1}{Z(x)}\,\piref(y\mid x)\exp\Bigl(\frac{r(x,y)}{\beta}\Bigr),\qquad Z(x)=\sum_{y}\piref(y\mid x)\,e^{r(x,y)/\beta},
\]
y el valor máximo es $\beta\log Z(x)$.
\end{teorema}

\addtocounter{definicion}{1}%

\begin{table}[h]
\centering\small
\begin{tabular}{@{}lccc@{}}
\toprule
$\beta$ & $\pi^{*}$ & $\E_{\pi^{*}}[r]$ & $\KL(\pi^{*}\,\|\,\piref)$ \\
\midrule
$\infty$ (referencia) & $(0{,}500;\ 0{,}300;\ 0{,}200)$ & $1{,}100$ & $0$ \\
$2$ & $(0{,}448;\ 0{,}443;\ 0{,}109)$ & $1{,}335$ & $0{,}058$ \\
$1$ & $(0{,}360;\ 0{,}587;\ 0{,}053)$ & $1{,}534$ & $0{,}205$ \\
$0{,}5$ & $(0{,}182;\ 0{,}808;\ 0{,}010)$ & $1{,}798$ & $0{,}587$ \\
\bottomrule
\end{tabular}
\caption{Política óptima del \cref{ej:beta} en función de $\beta$.}
\label{tab:beta}
\end{table}

\begin{observacion}[Distribución de Boltzmann y energía libre]\label{obs:boltzmann}
Con $E(y)=-r(y)$ y temperatura $\beta$, la política óptima es $\pi^{*}\propto\piref\,e^{-E/\beta}$: una distribución de Boltzmann (la softmax con temperatura del módulo~1) con $\piref$ como medida de fondo. El objetivo cambiado de signo, $\E_\pi[E]+\beta\,\KL(\pi\,\|\,\piref)$, es una energía libre, y la demostración del \cref{teo:optimo-kl} es el principio variacional de Gibbs: el mínimo vale $-\beta\log Z$. Cuando $\beta\to 0$, $\pi^{*}$ se concentra en las respuestas de máxima recompensa; cuando $\beta\to\infty$, $\pi^{*}\to\piref$ (\cref{eje:3-limites}).
\end{observacion}

El \cref{teo:optimo-kl} no puede aplicarse directamente a un LLM, porque $Z(x)$ es una suma sobre las $\abs{\V}^{T}$ respuestas (\cref{obs:bandido}). Por eso RLHF optimiza \eqref{eq:obj-kl} con PPO. El \cref{tema:5} muestra que la forma cerrada sí puede aprovecharse si se trabaja con diferencias de recompensas, en las que $Z(x)$ se cancela.


% =====================================================================
\chapter{Preferencias, modelo de recompensa y RLHF}
\label{tema:4}
% =====================================================================

Las personas no asignan con fiabilidad una nota absoluta a una respuesta, pero sí deciden con bastante consistencia cuál de dos respuestas prefieren. Este tema presenta el modelo de comparaciones por pares de Bradley y Terry~\cite{bt}, su estimación por máxima verosimilitud y su uso para entrenar un modelo de recompensa a partir de preferencias humanas, como en Christiano \emph{et al.}~\cite{christiano}, Stiennon \emph{et al.}~\cite{stiennon} y Ouyang \emph{et al.}~\cite{instructgpt}. Termina con el procedimiento completo de RLHF de InstructGPT. La parte de aprendizaje de recompensas sigue además la tarea~3 de CS234~\cite{cs234a3}.

\section{El modelo de Bradley--Terry}

\begin{definicion}[Modelo de Bradley--Terry]\label{def:bt}
Sean $n$ objetos con puntuaciones latentes $r_1,\dots,r_n\in\R$. El \emph{modelo de Bradley--Terry} postula que, al comparar $i$ con $j$, se prefiere $i$ con probabilidad
\[
P(i\pref j)=\frac{e^{r_i}}{e^{r_i}+e^{r_j}}=\sigma(r_i-r_j),
\]
y que las comparaciones son independientes entre sí.
\end{definicion}

La segunda igualdad se obtiene dividiendo numerador y denominador entre $e^{r_i}$. El modelo es una regresión logística (módulo~8) cuyo predictor es la \emph{diferencia} de puntuaciones.

\begin{proposicion}[Propiedades]\label{prop:bt}
\begin{enumerate}
  \item $P(i\pref j)+P(j\pref i)=1$, y $P(i\pref j)=\tfrac12$ si y solo si $r_i=r_j$.
  \item Las probabilidades no cambian si se suma una misma constante $c$ a todas las puntuaciones; las puntuaciones solo están determinadas salvo una traslación.
  \item Las \emph{razones de probabilidades} se multiplican: $\dfrac{P(i\pref j)}{P(j\pref i)}\cdot\dfrac{P(j\pref k)}{P(k\pref j)}=\dfrac{P(i\pref k)}{P(k\pref i)}$.
\end{enumerate}
\end{proposicion}

La propiedad (iii) es una forma de transitividad: si $i$ gana a $j$ con razón $3{:}1$ y $j$ a $k$ con razón $2{:}1$, el modelo predice que $i$ gana a $k$ con razón $6{:}1$, es decir, con probabilidad $6/7$.

\begin{definicion}[Modelo de Plackett--Luce]\label{def:pl}
Para ordenar $K$ objetos, el modelo de Plackett--Luce elige el primero con probabilidad proporcional a $e^{r_i}$, el segundo entre los restantes con la misma regla, y así sucesivamente:
\[
P(i_1\pref i_2\pref\dots\pref i_K)=\prod_{k=1}^{K}\frac{e^{r_{i_k}}}{\sum_{m\ge k}e^{r_{i_m}}}.
\]
Con $K=2$ se reduce al modelo de Bradley--Terry.
\end{definicion}

\section{Estimación por máxima verosimilitud}

\begin{proposicion}[Pérdida del modelo de recompensa]\label{prop:bt-mv}
Sea $\mathcal{D}=\{(x^{(m)},y_w^{(m)},y_l^{(m)})\}_{m=1}^{M}$ un conjunto de comparaciones en las que $y_w$ fue preferida a $y_l$, y sea $r_\phi(x,y)$ una recompensa parametrizada. Bajo el modelo de Bradley--Terry, el estimador de máxima verosimilitud de $\phi$ minimiza
\[
\LRM(\phi)=-\frac1M\sum_{m=1}^{M}\log\sigma\bigl(r_\phi(x^{(m)},y_w^{(m)})-r_\phi(x^{(m)},y_l^{(m)})\bigr).
\]
Con $z=r_\phi(x,y_w)-r_\phi(x,y_l)$, el gradiente de cada término es $-(1-\sigma(z))\bigl(\nabla_\phi r_\phi(x,y_w)-\nabla_\phi r_\phi(x,y_l)\bigr)$.
\end{proposicion}

El factor $1-\sigma(z)=\sigma(-z)$ es la probabilidad que el modelo asigna al orden \emph{equivocado}: los pares ya bien ordenados apenas contribuyen y los mal ordenados empujan con fuerza.

\addtocounter{definicion}{1}%

\begin{proposicion}[Estimador con dos objetos]\label{prop:bt-dos}
Si en $n$ comparaciones entre $a$ y $b$ se prefiere $a$ en $n_a$ ocasiones, con $0<n_a<n$, el estimador de máxima verosimilitud de la diferencia es $\widehat{r_a-r_b}=\log\bigl(n_a/(n-n_a)\bigr)$, y la probabilidad estimada es la frecuencia observada $n_a/n$.
\end{proposicion}

Si $n_a=n$, la verosimilitud crece sin cota cuando $z\to\infty$ y el estimador no existe: un modelo de recompensa sobreajustado puede separar las preferencias de entrenamiento con márgenes arbitrariamente grandes. En la práctica lo impiden la regularización y la parada temprana.

\section{Modelos de recompensa neuronales}

En el ajuste de LLM, $r_\phi(x,y)$ es un Transformer, normalmente inicializado desde el modelo SFT, en el que la capa de salida sobre el vocabulario se sustituye por una cabeza lineal que produce un escalar a partir del estado oculto del último token~\cite{stiennon,instructgpt}. Tres detalles merecen mención.
\begin{itemize}[leftmargin=1.4em,label=\textendash]
  \item \emph{Comparaciones de segmentos.} Christiano \emph{et al.}~\cite{christiano} aplicaron Bradley--Terry a pares de fragmentos de trayectoria en tareas de control, con la recompensa de un fragmento igual a la suma de recompensas por paso predichas (\cref{eje:4-cs234}).
  \item \emph{Ordenaciones.} En InstructGPT, los anotadores ordenaban entre $K=4$ y $K=9$ respuestas a la misma instrucción. Cada ordenación produce $\binom{K}{2}$ comparaciones, que se tratan en un mismo lote y se ponderan por $1/\binom{K}{2}$ para que las instrucciones con más respuestas no dominen~\cite[§3.5]{instructgpt}.
  \item \emph{Normalización.} Por la \cref{prop:bt}(ii), la recompensa aprendida está definida salvo una constante; se suele trasladar para que las demostraciones de referencia tengan recompensa media nula antes del RL~\cite{stiennon,instructgpt}.
\end{itemize}

\section{El procedimiento RLHF de InstructGPT}

\addtocounter{figure}{1}%

El procedimiento de Ouyang \emph{et al.}~\cite[§3.5]{instructgpt} consta de tres etapas (\cref{fig:rlhf}):
\begin{enumerate}[label=\arabic*.]
  \item \emph{Ajuste supervisado}: el modelo preentrenado se ajusta con demostraciones escritas por anotadores (módulo~20).
  \item \emph{Modelo de recompensa}: se recogen ordenaciones de respuestas del modelo SFT y se entrena $r_\phi$ con la pérdida de la \cref{prop:bt-mv}. En el artículo, los modelos de recompensa tienen $6\cdot 10^{9}$ parámetros.
  \item \emph{Optimización por refuerzo}: se maximiza $\E[r_\phi(x,y)]-\beta\,\KL(\pi_\theta\,\|\,\piref)$ con PPO (\cref{tema:3}). La variante PPO-ptx añade un término $\gamma\,\E_{x\sim\mathcal{D}_{\mathrm{pre}}}[\log\pi_\theta(x)]$ de verosimilitud sobre datos de preentrenamiento para reducir la pérdida de rendimiento en tareas clásicas.
\end{enumerate}

Con este procedimiento, las respuestas de un modelo de $1{,}3\cdot 10^{9}$ parámetros fueron preferidas por los evaluadores a las de GPT-3 de $1{,}75\cdot 10^{11}$ parámetros~\cite{instructgpt}. El coste práctico es alto: durante la tercera etapa conviven en memoria la política, la referencia, el modelo de recompensa y el crítico de PPO, cuatro redes del tamaño de un LLM. Esta es la principal motivación práctica de DPO (\cref{tema:5}).


% =====================================================================
\chapter{Optimización directa de preferencias (DPO)}
\label{tema:5}
% =====================================================================

La optimización directa de preferencias (DPO) de Rafailov \emph{et al.}~\cite{dpo} elimina el modelo de recompensa explícito y el bucle de refuerzo de RLHF. La idea es invertir la forma cerrada del \cref{teo:optimo-kl} para expresar la recompensa en función de la política y sustituirla en el modelo de Bradley--Terry, donde la función de partición se cancela. Este tema reproduce la derivación completa, los resultados de equivalencia de la §5.1 del artículo y el gradiente de la §4. El material se complementa con la clase invitada sobre DPO del curso CS234~\cite{cs234} y con la tarea~5 de CS336~\cite{cs336a5s}.

\section{De la política óptima a la recompensa}

\begin{lema}[Reparametrización]\label{lem:reparam}
Sea $\pi_r^{*}$ la política óptima del objetivo \eqref{eq:obj-kl} para la recompensa $r$ y la referencia $\piref$, y supóngase $\piref(y\mid x)>0$. Entonces, para todo $(x,y)$,
\[
r(x,y)=\beta\log\frac{\pi_r^{*}(y\mid x)}{\piref(y\mid x)}+\beta\log Z(x).
\]
\end{lema}

El término incómodo, $\beta\log Z(x)$, depende solo de la instrucción $x$ y no de la respuesta. Todo lo que sigue explota este hecho.

\section{Clases de equivalencia de recompensas}

\begin{definicion}[Recompensas equivalentes]\label{def:equiv}
Dos recompensas $r$ y $r'$ son \emph{equivalentes} si existe una función $f$ tal que $r'(x,y)=r(x,y)+f(x)$ para todo $(x,y)$. Es una relación de equivalencia; sus clases se llaman \emph{clases de recompensas}.
\end{definicion}

\begin{lema}[Las recompensas equivalentes inducen las mismas preferencias]\label{lem:equiv-pref}
Si $r$ y $r'$ son equivalentes, inducen la misma distribución de preferencias de Bradley--Terry, y también de Plackett--Luce, sobre las respuestas a cada instrucción.
\end{lema}

\begin{lema}[Las recompensas equivalentes inducen la misma política óptima]\label{lem:equiv-pol}
Si $r$ y $r'$ son equivalentes, $\pi_{r'}^{*}=\pi_r^{*}$.
\end{lema}

\begin{teorema}[Representación de las clases de recompensas]\label{teo:representacion}
Toda clase de recompensas contiene un representante de la forma $\hat r(x,y)=\beta\log\bigl(\pi(y\mid x)/\piref(y\mid x)\bigr)$ para alguna política $\pi$, y esa política es única: es la política óptima común de la clase.
\end{teorema}

El teorema~1 de~\cite{dpo} enuncia este resultado: parametrizar las recompensas como $\beta\log(\pi/\piref)$ no pierde generalidad, porque toda clase de recompensas compatible con el modelo de preferencias puede escribirse así. Es el sentido preciso del subtítulo del artículo: el modelo de lenguaje es, en secreto, un modelo de recompensa.

\section{La pérdida DPO}

\begin{proposicion}[Preferencias en función de la política óptima]\label{prop:pref-pol}
Si las preferencias siguen el modelo de Bradley--Terry con recompensa $r$ y $\pi^{*}=\pi^{*}_r$, entonces
\[
P(y_w\pref y_l\mid x)=\sigma\Bigl(\beta\log\frac{\pi^{*}(y_w\mid x)}{\piref(y_w\mid x)}-\beta\log\frac{\pi^{*}(y_l\mid x)}{\piref(y_l\mid x)}\Bigr).
\]
\end{proposicion}

\begin{definicion}[Pérdida DPO]\label{def:dpo}
Sea $\pi_\theta$ una política parametrizada y $\hat r_\theta(x,y)=\beta\log\bigl(\pi_\theta(y\mid x)/\piref(y\mid x)\bigr)$ su \emph{recompensa implícita}. La pérdida DPO sobre un conjunto de comparaciones $\mathcal{D}$ es
\[
\LDPO(\theta)=-\E_{(x,y_w,y_l)\sim\mathcal{D}}\Bigl[\log\sigma\Bigl(\beta\log\frac{\pi_\theta(y_w\mid x)}{\piref(y_w\mid x)}-\beta\log\frac{\pi_\theta(y_l\mid x)}{\piref(y_l\mid x)}\Bigr)\Bigr].
\]
\end{definicion}

La derivación completa se resume en cuatro pasos: (1) el objetivo con KL tiene la solución cerrada del \cref{teo:optimo-kl}; (2) se despeja la recompensa (\cref{lem:reparam}); (3) se sustituye en Bradley--Terry y se cancela $Z(x)$ (\cref{prop:pref-pol}); (4) se reemplaza $\pi^{*}$ por $\pi_\theta$ y se maximiza la verosimilitud de las preferencias. El paso~(4) es legítimo por el \cref{teo:representacion}: ajustar $\pi_\theta$ equivale a ajustar una recompensa dentro de una familia que, sin restricciones de capacidad, contiene todas las clases. Con otras palabras, $\LDPO$ es la pérdida $\LRM$ de la \cref{prop:bt-mv} con $r_\phi$ sustituida por $\hat r_\theta$, y si $\hat r_\theta$ minimiza $\LRM$, entonces $\pi_\theta$ es la política óptima del objetivo con KL para esa recompensa.

En la implementación, $\log\pi_\theta(y\mid x)=\sum_t\log\pi_\theta(y_t\mid x,y_{<t})$ se obtiene sumando las log-probabilidades por token de la respuesta (\cref{prop:mdp-tokens}), y las de $\piref$ se calculan una vez con el modelo congelado.

\begin{proposicion}[Gradiente de DPO]\label{prop:grad-dpo}
Sea $h=\hat r_\theta(x,y_w)-\hat r_\theta(x,y_l)$. Entonces
\[
\nabla_\theta\LDPO=-\beta\,\E_{\mathcal{D}}\Bigl[\sigma(-h)\bigl(\nabla_\theta\log\pi_\theta(y_w\mid x)-\nabla_\theta\log\pi_\theta(y_l\mid x)\bigr)\Bigr].
\]
\end{proposicion}

El gradiente aumenta la log-probabilidad de la respuesta preferida y reduce la de la rechazada. El peso $\sigma(-h)$ es la probabilidad que la recompensa implícita asigna al orden equivocado: es grande cuando el modelo ordena mal el par y casi nulo cuando ya lo ordena bien. Rafailov \emph{et al.} observan que sin este peso el método degenera~\cite[§4]{dpo}.

\addtocounter{definicion}{1}%

\section{DPO frente a RLHF con PPO}

\begin{table}[h]
\centering\small
\begin{tabularx}{\textwidth}{@{}l>{\raggedright\arraybackslash}X>{\raggedright\arraybackslash}X@{}}
\toprule
 & RLHF con PPO & DPO \\
\midrule
Redes en memoria & 4 (política, referencia, recompensa, crítico) & 2 (política, referencia) \\
Muestreo al entrenar & sí (\emph{on-policy}) & no (datos fijos, \emph{offline}) \\
Pérdida & objetivo recortado y penalización KL & clasificación binaria \\
Óptimo teórico & $\pi^{*}$ del objetivo con KL & $\pi^{*}$ del objetivo con KL \\
\bottomrule
\end{tabularx}
\caption{Comparación entre RLHF con PPO y DPO.}
\label{tab:dpo-ppo}
\end{table}

Ambos métodos persiguen el mismo óptimo (\cref{tab:dpo-ppo}), pero no son equivalentes en la práctica. DPO solo ve las respuestas del conjunto de datos, generadas en general por otra política, y no explora; PPO evalúa con el modelo de recompensa las respuestas que la política genera en cada momento. Además, DPO controla solo la diferencia $h$: puede reducir a la vez las probabilidades de ambas respuestas siempre que la rechazada baje más, y puede sobreoptimizar su recompensa implícita, por ejemplo favoreciendo respuestas más largas (\cref{tema:6}). El apéndice~A.3 de~\cite{dpo} extiende la pérdida a ordenaciones completas con el modelo de Plackett--Luce.


% =====================================================================
\chapter{Retroalimentación de IA, recompensas verificables y sobreoptimización}
\label{tema:6}
% =====================================================================

Los temas anteriores suponen que las preferencias proceden de personas y que la recompensa es fiable. Este tema estudia tres desarrollos que cuestionan esas hipótesis: la sustitución de parte de la supervisión humana por retroalimentación generada por un modelo según una lista de principios, Constitutional AI~\cite{cai}; el uso de verificadores automáticos como recompensa, con el algoritmo GRPO~\cite{deepseekmath,r1} y su análisis en la clase~17 y la tarea~5 de CS336~\cite{cs336,cs336a5}, y la sobreoptimización de recompensas imperfectas, medida por Gao, Schulman y Hilton~\cite{gao}.

\section{Constitutional AI y retroalimentación de IA}

Constitutional AI (Bai \emph{et al.}, Anthropic, 2022)~\cite{cai} entrena un asistente inofensivo sin etiquetas humanas sobre inocuidad: la supervisión humana se concentra en redactar una lista de principios en lenguaje natural, la \emph{constitución}. El método tiene dos fases (\cref{fig:cai}).
\begin{enumerate}[label=\arabic*.]
  \item \emph{Fase supervisada (SL-CAI).} Un modelo útil responde a instrucciones problemáticas. Después se le pide criticar su propia respuesta según un principio elegido al azar y reescribirla; el proceso puede repetirse. El modelo se ajusta con SFT sobre las revisiones finales.
  \item \emph{Fase de refuerzo (RL-CAI).} El modelo de la fase~1 genera pares de respuestas. Un modelo de retroalimentación decide cuál cumple mejor un principio, en algunos experimentos razonando paso a paso. Con esas preferencias generadas por IA se entrena un modelo de preferencias de Bradley--Terry (\cref{tema:4}) y se optimiza la política contra él con RL y penalización KL (\cref{tema:3}). Es lo que se denomina RLAIF, \emph{RL from AI feedback}.
\end{enumerate}

\addtocounter{figure}{1}%

Las matemáticas son las de los temas~3 y~4; lo que cambia es quién genera las comparaciones. En el artículo, las preferencias de utilidad siguen procediendo de personas y las de inocuidad, del modelo~\cite{cai}. La ventaja es doble: el método escala con mucha menos anotación humana, y los valores que se optimizan quedan explícitos y pueden auditarse. Anthropic ha descrito públicamente Constitutional AI como parte del entrenamiento de Claude; la constitución concreta de cada versión y los detalles actuales de su entrenamiento solo se conocen en la medida en que se han publicado.

\section{Recompensas verificables}

En matemáticas o programación no hace falta aprender la recompensa: se puede comprobar la respuesta.

\begin{definicion}[Recompensa verificable]\label{def:rlvr}
Una \emph{recompensa verificable} es una función $r(x,y)\in\{0,1\}$ calculada por un programa: vale~1 si la respuesta final extraída de $y$ coincide con la solución de referencia o si el código de $y$ supera unas pruebas, y~0 en otro caso. El RL con recompensas de este tipo se denomina RLVR (\emph{reinforcement learning with verifiable rewards}), término introducido en Tülu~3~\cite{tulu3}.
\end{definicion}

Como solo se premia el resultado, el modelo es libre de escribir cualquier texto intermedio. DeepSeek-R1~\cite{r1} documenta que, al optimizar recompensas de exactitud y de formato con GRPO directamente sobre un modelo base, sin SFT previo (el modelo R1-Zero), las respuestas se alargan y aparecen cadenas de razonamiento extensas que aumentan la tasa de acierto. Sobre los modelos cerrados la información pública es escasa: OpenAI declaró que o1 (2024) se entrenó con RL a gran escala para razonar mediante cadenas de pensamiento sin publicar el algoritmo; Anthropic ofrece un modo de pensamiento extendido en Claude desde Claude~3.7 Sonnet (2025), y Google ofrece modos de razonamiento en Gemini. Los detalles de su entrenamiento no son públicos.

\section{GRPO: optimización de política relativa al grupo}

PPO necesita un crítico del tamaño del modelo para estimar $V(s_t)$. GRPO, introducido en DeepSeekMath~\cite[§4.1]{deepseekmath}, lo sustituye por una línea base empírica: para cada instrucción se generan varias respuestas y cada una se compara con las demás.

\begin{definicion}[Ventaja relativa al grupo]\label{def:grpo-ventaja}
Para una instrucción $x$ se muestrean $G$ respuestas $y^{(1)},\dots,y^{(G)}\sim\piold(\cdot\mid x)$ con recompensas $r_1,\dots,r_G$. La ventaja de la respuesta $i$, asignada a todos sus tokens, es
\[
\widehat A_i=\frac{r_i-\media(r_1,\dots,r_G)}{\desv(r_1,\dots,r_G)+\varepsilon_0},
\]
con $\varepsilon_0>0$ pequeño para evitar la división por cero.
\end{definicion}

\begin{definicion}[Objetivo de GRPO]\label{def:grpo}
Con $\rho_{i,t}=\pi_\theta(y^{(i)}_t\mid x,y^{(i)}_{<t})/\piold(y^{(i)}_t\mid x,y^{(i)}_{<t})$, el objetivo que se maximiza es~\cite[ec.~(3)]{deepseekmath}
\[
\frac1G\sum_{i=1}^{G}\frac{1}{\abs{y^{(i)}}}\sum_{t=1}^{\abs{y^{(i)}}}\Bigl[\min\bigl(\rho_{i,t}\widehat A_i,\ \clip(\rho_{i,t},1-\varepsilon,1+\varepsilon)\widehat A_i\bigr)-\beta\,\widehat{\KL}_{i,t}\Bigr],
\]
donde $\widehat{\KL}_{i,t}=u-\log u-1$ con $u=\piref(y^{(i)}_t\mid\cdot)/\pi_\theta(y^{(i)}_t\mid\cdot)$.
\end{definicion}

A diferencia de RLHF con PPO, la penalización KL no se resta de la recompensa sino que se añade directamente a la pérdida, y se estima con el estimador $k_3$ del \cref{eje:3-estimadores}.

\begin{proposicion}[Propiedades del estimador $k_3$]\label{prop:k3}
Sean $p$ y $q$ distribuciones con $q(y)>0$ donde $p(y)>0$, y sea $k_3(y)=u-\log u-1$ con $u=q(y)/p(y)$. Entonces $k_3(y)\ge 0$ para todo $y$ y $\E_{y\sim p}[k_3(y)]=\KL(p\,\|\,q)$.
\end{proposicion}

\begin{proposicion}[Propiedades de las ventajas de grupo]\label{prop:grpo}
Con la desviación poblacional y $\varepsilon_0=0$, si no todas las recompensas son iguales, las ventajas de un grupo tienen media~0 y varianza~1. Si todas las recompensas son iguales, todas las ventajas son nulas y el grupo no aporta gradiente al término recortado.
\end{proposicion}

La media del grupo desempeña el papel de la línea base del \cref{tema:2}. Con recompensas binarias y probabilidad de acierto $p$, un grupo de $G$ respuestas no aporta señal con probabilidad $p^{G}+(1-p)^{G}$: las instrucciones que el modelo casi siempre resuelve, o casi nunca, desperdician cómputo (\cref{eje:6-cero}).

\addtocounter{definicion}{1}%

\begin{observacion}[Sesgos de la normalización]\label{obs:drgrpo}
Liu \emph{et al.}~\cite{drgrpo} señalan dos sesgos en el objetivo de la \cref{def:grpo}. Dividir por la desviación del grupo da más peso a las instrucciones con poca variación de recompensa, y dividir por la longitud $\abs{y^{(i)}}$ hace que, con ventaja negativa, cada token de una respuesta larga se penalice menos que uno de una corta, lo que favorece respuestas incorrectas largas. Su variante, Dr.~GRPO, elimina ambas normalizaciones (\cref{eje:6-drgrpo,eje:6-longitud}).
\end{observacion}

\section{Sobreoptimización de la recompensa}

Cuando se optimiza un indicador imperfecto, la política acaba encontrando formas de aumentarlo sin mejorar lo que el indicador pretendía medir. Es la ley de Goodhart aplicada al refuerzo, conocida como \emph{reward hacking} (\cref{tab:hacking}).

\begin{table}[h]
\centering\small
\begin{tabular}{@{}ll@{}}
\toprule
Indicador & Explotación típica \\
\midrule
modelo de recompensa de utilidad & respuestas más largas, más listas y más halagos \\
pruebas unitarias & código que detecta la prueba y devuelve el valor esperado \\
satisfacción del usuario & dar la razón al usuario (adulación, \emph{sycophancy}) \\
formato correcto & etiquetas impecables con contenido vacío \\
\bottomrule
\end{tabular}
\caption{Ejemplos de explotación de recompensas.}
\label{tab:hacking}
\end{table}

Gao, Schulman y Hilton~\cite{gao} midieron el fenómeno de forma controlada: un modelo de recompensa grande hace de recompensa «verdadera» y etiqueta los datos con los que se entrena un modelo de recompensa \emph{proxy} más pequeño; después se optimiza la política contra el proxy y se mide la recompensa verdadera. Al aumentar la divergencia respecto a la política inicial, la recompensa proxy sigue creciendo, mientras que la verdadera crece, alcanza un máximo y decrece.

\begin{definicion}[Mejor de $n$]\label{def:bon}
La política \emph{mejor de $n$} (\emph{best-of-$n$}) genera $n$ respuestas independientes con la política inicial y devuelve la de mayor recompensa proxy.
\end{definicion}

Con $d=\sqrt{\KL(\pi\,\|\,\pi_{\mathrm{init}})}$, Gao \emph{et al.} ajustan las formas funcionales
\[
R_{\mathrm{bon}}(d)=d\,(\alpha_{\mathrm{bon}}-\beta_{\mathrm{bon}}\,d),\qquad R_{\mathrm{RL}}(d)=d\,(\alpha_{\mathrm{RL}}-\beta_{\mathrm{RL}}\log d)
\]
para la ganancia de recompensa verdadera con mejor de $n$ y con RL, con coeficientes que varían de forma suave con el tamaño del modelo de recompensa~\cite{gao}. Para mejor de $n$, la divergencia se calcula analíticamente: $\mathrm{KL}_{\mathrm{bon}}=\log n-(n-1)/n$ (\cref{eje:6-bon}).

\begin{proposicion}[Divergencia óptima]\label{prop:gao}
$R_{\mathrm{bon}}$ alcanza su máximo en $d^{*}=\alpha_{\mathrm{bon}}/(2\beta_{\mathrm{bon}})$, con valor $\alpha_{\mathrm{bon}}^{2}/(4\beta_{\mathrm{bon}})$. $R_{\mathrm{RL}}$ alcanza su máximo en $d^{*}=e^{\alpha_{\mathrm{RL}}/\beta_{\mathrm{RL}}-1}$, con valor $\beta_{\mathrm{RL}}d^{*}$.
\end{proposicion}

Existe, por tanto, una cantidad óptima de optimización, y más allá de ella la política empeora. Este es el argumento cuantitativo a favor de la penalización KL del \cref{tema:3}: limita cuánto puede explotar la política los errores de $r_\phi$. Otras defensas habituales son la parada temprana, el reentrenamiento del modelo de recompensa con muestras de la política actual, el uso de recompensas verificables cuando existen, la combinación de varios modelos de recompensa y la revisión manual de las muestras con recompensa anómalamente alta. DPO no es inmune: también puede sobreoptimizar su recompensa implícita.



\backmatter
% =====================================================================
\begin{thebibliography}{99}
\addcontentsline{toc}{chapter}{Bibliografía}
\markboth{Bibliografía}{}
\small

\bibitem{cai}
Y.~Bai, S.~Kadavath, S.~Kundu, A.~Askell \emph{et al.}
\newblock Constitutional AI: Harmlessness from AI feedback.
\newblock Prepublicación, Anthropic, 2022.
\newblock \url{https://arxiv.org/abs/2212.08073}

\bibitem{bt}
R.~A. Bradley y M.~E. Terry.
\newblock Rank analysis of incomplete block designs: I. The method of paired comparisons.
\newblock \emph{Biometrika}, 39(3/4):324--345, 1952.

\bibitem{christiano}
P.~Christiano, J.~Leike, T.~B. Brown, M.~Martic, S.~Legg y D.~Amodei.
\newblock Deep reinforcement learning from human preferences.
\newblock En \emph{Advances in Neural Information Processing Systems (NeurIPS)}, 2017.
\newblock \url{https://arxiv.org/abs/1706.03741}

\bibitem{r1}
DeepSeek-AI (D.~Guo, D.~Yang, H.~Zhang, J.~Song \emph{et al.}).
\newblock DeepSeek-R1: Incentivizing reasoning capability in LLMs via reinforcement learning.
\newblock Prepublicación, 2025. Publicado en \emph{Nature}, 645:633--638, 2025.
\newblock \url{https://arxiv.org/abs/2501.12948}

\bibitem{cs234a3}
E.~Brunskill y equipo docente de CS234.
\newblock CS234: Reinforcement Learning. Tarea~3.
\newblock Stanford University, invierno de 2025. Secciones~1 «Reward engineering», 2 «Learning from preferences» y~3 «Direct preference optimization».
\newblock \url{https://web.stanford.edu/class/cs234/assignments/a3/hw3_questions.pdf}

\bibitem{cs234}
E.~Brunskill.
\newblock CS234: Reinforcement Learning.
\newblock Stanford University, primavera de 2024. Bloque «Policy Gradient» (clases~5--7) y clase~9, con la clase invitada sobre DPO de R.~Rafailov, A.~Sharma y E.~Mitchell.
\newblock \url{https://web.stanford.edu/class/cs234/CS234Spr2024/modules.html}

\bibitem{cs336}
P.~Liang, T.~Hashimoto \emph{et al.}
\newblock CS336: Language Modeling from Scratch.
\newblock Stanford University, primavera de 2025. Clase~17: aprendizaje por refuerzo con recompensas verificables, gradiente de política y GRPO.
\newblock \url{https://github.com/stanford-cs336/spring2025-lectures}

\bibitem{cs336a5}
Equipo docente de CS336.
\newblock CS336: Language Modeling from Scratch. Tarea~5, «Alignment».
\newblock Stanford University, primavera de 2025. §7.2 y §8.
\newblock \url{https://github.com/stanford-cs336/assignment5-alignment}

\bibitem{cs336a5s}
Equipo docente de CS336.
\newblock CS336: Language Modeling from Scratch. Suplemento de la tarea~5, «Instruction Tuning and RLHF».
\newblock Stanford University, primavera de 2026. §6.3.
\newblock \url{https://github.com/stanford-cs336/assignment5-alignment}

\bibitem{gao}
L.~Gao, J.~Schulman y J.~Hilton.
\newblock Scaling laws for reward model overoptimization.
\newblock En \emph{International Conference on Machine Learning (ICML)}, 2023.
\newblock \url{https://arxiv.org/abs/2210.10760}

\bibitem{tulu3}
N.~Lambert, J.~Morrison, V.~Pyatkin, S.~Huang, H.~Ivison \emph{et al.}
\newblock Tülu 3: Pushing frontiers in open language model post-training.
\newblock Prepublicación, Allen Institute for AI, 2024.
\newblock \url{https://arxiv.org/abs/2411.15124}

\bibitem{cs285}
S.~Levine.
\newblock CS285: Deep Reinforcement Learning.
\newblock University of California, Berkeley, otoño de 2023. Clases~4 «Introduction to Reinforcement Learning», 5 «Policy Gradients», 6 «Actor-Critic Algorithms» y~9 «Advanced Policy Gradients».
\newblock \url{https://rail.eecs.berkeley.edu/deeprlcourse-fa23/}

\bibitem{cs285hw2}
S.~Levine y equipo docente de CS285.
\newblock CS285: Deep Reinforcement Learning. Tarea~2, «Policy Gradients».
\newblock University of California, Berkeley, otoño de 2023. §2 (repaso) y §8 (problemas de análisis).
\newblock \url{https://rail.eecs.berkeley.edu/deeprlcourse-fa23/}

\bibitem{drgrpo}
Z.~Liu, C.~Chen, W.~Li, P.~Qi, T.~Pang, C.~Du, W.~S. Lee y M.~Lin.
\newblock Understanding R1-Zero-like training: A critical perspective.
\newblock Prepublicación, 2025.
\newblock \url{https://arxiv.org/abs/2503.20783}

\bibitem{instructgpt}
L.~Ouyang, J.~Wu, X.~Jiang, D.~Almeida, C.~L. Wainwright \emph{et al.}
\newblock Training language models to follow instructions with human feedback.
\newblock En \emph{Advances in Neural Information Processing Systems (NeurIPS)}, 2022.
\newblock \url{https://arxiv.org/abs/2203.02155}

\bibitem{dpo}
R.~Rafailov, A.~Sharma, E.~Mitchell, S.~Ermon, C.~D. Manning y C.~Finn.
\newblock Direct preference optimization: Your language model is secretly a reward model.
\newblock En \emph{Advances in Neural Information Processing Systems (NeurIPS)}, 2023.
\newblock \url{https://arxiv.org/abs/2305.18290}

\bibitem{gae}
J.~Schulman, P.~Moritz, S.~Levine, M.~I. Jordan y P.~Abbeel.
\newblock High-dimensional continuous control using generalized advantage estimation.
\newblock En \emph{International Conference on Learning Representations (ICLR)}, 2016.
\newblock \url{https://arxiv.org/abs/1506.02438}

\bibitem{ppo}
J.~Schulman, F.~Wolski, P.~Dhariwal, A.~Radford y O.~Klimov.
\newblock Proximal policy optimization algorithms.
\newblock Prepublicación, 2017.
\newblock \url{https://arxiv.org/abs/1707.06347}

\bibitem{deepseekmath}
Z.~Shao, P.~Wang, Q.~Zhu, R.~Xu, J.~Song \emph{et al.}
\newblock DeepSeekMath: Pushing the limits of mathematical reasoning in open language models.
\newblock Prepublicación, 2024.
\newblock \url{https://arxiv.org/abs/2402.03300}

\bibitem{silver}
D.~Silver.
\newblock UCL Course on Reinforcement Learning.
\newblock University College London, 2015. Clase~2, «Markov Decision Processes».
\newblock \url{https://www.davidsilver.uk/teaching/}

\bibitem{stiennon}
N.~Stiennon, L.~Ouyang, J.~Wu, D.~M. Ziegler, R.~Lowe, C.~Voss, A.~Radford, D.~Amodei y P.~Christiano.
\newblock Learning to summarize from human feedback.
\newblock En \emph{Advances in Neural Information Processing Systems (NeurIPS)}, 2020.
\newblock \url{https://arxiv.org/abs/2009.01325}

\bibitem{sb}
R.~S. Sutton y A.~G. Barto.
\newblock \emph{Reinforcement Learning: An Introduction}, 2.ª ed.
\newblock MIT Press, 2018. Caps.~3 «Finite Markov Decision Processes» y~13 «Policy Gradient Methods».
\newblock \url{http://incompleteideas.net/book/the-book-2nd.html}

\bibitem{williams}
R.~J. Williams.
\newblock Simple statistical gradient-following algorithms for connectionist reinforcement learning.
\newblock \emph{Machine Learning}, 8(3--4):229--256, 1992.

\end{thebibliography}


\end{document}
